\section{Few-Shot Image Generation in GANs}

Few-shot image generation (FSIG) extends few-shot learning to generative modeling: 
the goal is to synthesize novel images for a new category given only a handful of training examples. 
This contrasts with few-shot classification, where techniques like transfer learning have proven effective \citep{chowdhury2021fewshotimageclassificationjust}.
In FSIG, however, naively fine-tuning a powerful generative model on few examples leads to overfitting and mode collapse.
According to \citet{zhao2023closerlookfewshotimage}, modern GANs “tend to replicate the training samples” when adapted to only 10 target images.
As a result, FSIG methods typically transfer rich prior knowledge from a large source domain and then adapt carefully to the new domain. 
This is often framed as a style-transfer or domain-adaptation problem: 
a generator pretrained on a large dataset ,like FFHQ faces, is fine-tuned on a small target set without access to the source data \citep{cao2024fewshotimagegenerationconditional,zhao2023closerlookfewshotimage}.
The key challenge is preserving the source model’s diversity and high-fidelity output while learning the new style from very limited data.

Generative Adversarial Networks (GANs) have been the dominant framework for FSIG. 
The standard paradigm is to take a GAN (often StyleGAN or BigGAN) pretrained on a 
large source dataset and adapt it to the new class. Early approaches simply fine-tuned all weights on the few-shot 
target (TGAN) \citep{zhao2023closerlookfewshotimage}, but this quickly overfits (the generator memorizes the few examples). Numerous strategies have been proposed to prevent collapse and preserve diversity during adaptation.
or instance, FreezeD freezes some high-resolution layers of the discriminator, ADA (adaptive data augmentation) introduces random augmentations during training, and EWC applies an Elastic Weight Consolidation penalty on important weights \citep{mo2020freezediscriminatorsimplebaseline}.
Another approach Cross-domain Correspondence (CDC) preserves the relative distances between different source-domain samples in the target embedding space \citep{ojha2021fewshotimagegenerationcrossdomain}.
In practice, Zhao et al. (2022) found that most of these GAN-adaptation methods converge to similar image quality after many iterations – the real difference is how slowly they lose diversity during training \citep{zhao2023closerlookfewshotimage}.
In other words, diversity degradation (mode collapse) is the main bottleneck: methods that slow down this collapse (e.g. by freezing parameters or using strong regularization) tend to perform better overall.

Building on this insight, Zhao et al. propose a contrastive learning scheme to retain diversity. 
Their Dual Contrastive Learning (DCL) maximizes mutual information between the source and target generators: 
the idea is that the same latent input z should map to semantically corresponding images before and after adaptation.
Concretely, they use a contrastive loss on multi-level features of the generator and discriminator to enforce that 
source and target outputs remain similar in high-level content. This preserves the source generator's 
rich variation in the adapted model, yielding state-of-the-art results on several datasets \citep{zhao2023closerlookfewshotimage}.
Similarly, Hong et al. (2020) introduce DeltaGAN, which decomposes generation into a reconstruction sub-network 
(that computes an intra-class “delta” transformation between two samples of the same category) and a generation sub-network 
(that applies a sample-specific delta to a base image) \citep{hong2022deltagandiversefewshotimage}.
DeltaGAN’s adversarial “delta matching” encourages high diversity by learning many distinct transformations from the limited examples.

GAN-based FSIG methods share a transfer-learning theme: 
leverage a pretrained model's knowledge and carefully regularize the adaptation. 
Typical training strategies include fine-tuning, parameter freezing, data augmentation, and contrastive or distance-preserving losses. 
Their success is usually measured by metrics like Fréchet Inception Distance (FID) or LPIPS (perceptual diversity).

Across these models, several common training strategies emerge. Transfer Learning is most common: one typically starts from a pretrained generative model (GAN, VAE, or diffusion) 
and updates it using the few target images. Pure fine-tuning (e.g. updating all weights) generally overfits, so many methods constrain the update. 
Examples include freezing layers (FreezeD) or parameters (like BatchNorm scaling), domain-regularization, and data augmentation (ADA) to simulate more data \citep{zhao2023closerlookfewshotimage}.
Meta-Learning is another strategy: the model is trained on many simulated few-shot tasks so that it can adapt quickly. 
As noted, meta-learned GANs/VAEs (e.g. using MAML/Reptile) have been attempted, 
but often still need per-task fine-tuning and can struggle to generate sharp images \citep{hong2022deltagandiversefewshotimage}.
Contrastive/Distance Losses have become popular: enforcing that the adapted generator preserves relationships (distances or class labels) from the source domain helps maintain diversity.
Finally, data augmentation and synthesis can create pseudo-examples: e.g. methods like AdaGAN or mixing source images into the target fine-tuning process.

Evaluation strategies are also noteworthy. 
FSIG is evaluated on cross-domain tasks such as generating artistic or cartoon versions of FFHQ faces, 
or adapting LSUN Church to “Haunted House” style, or similar sketch-to-photo tasks. Metrics include FID 
for quality (when enough data exists) and LPIPS or recall for diversity. \citet{zhao2023closerlookfewshotimage} emphasize that with only 10 real 
images, FID is unstable, so they instead use intra-LPIPS (average pairwise LPIPS of generated set) to quantify diversity .

Despite these advances, few-shot image generation remains challenging. Surveys note that learning from very few or zero examples is still “a serious challenge” \citep{song2022comprehensivesurveyfewshotlearning}.
Key open issues include: catastrophic forgetting (the model loses the source distribution entirely), 
mode collapse (outputs lack diversity), and evaluation (how to measure quality/diversity with tiny real sample sets). 
Overfitting to the few examples is ubiquitous: even with strong regularization, generators tend to memorize salient features. 
Furthermore, different generative families have trade-offs – GANs are fast and high-quality but brittle, VAEs are stable but blurrier, 
and diffusion models are powerful but computationally heavy. A brief comparison: GANs excel at sharp images but require careful balancing of 
losses; VAEs naturally encourage covering the whole distribution (which aids diversity) at the cost of some sharpness; 
diffusion models sample iteratively and often capture fine detail, but naively adapting them can be even more parameter-sensitive than GANs \citep{cao2024fewshotimagegenerationconditional}.




